\section{The exact local Hilbert completion and its cutoff-transport obstruction}
\label{sec:local-hilbert-cutoff-obstruction}

\paragraph{Status and purpose.}
The complete post-writing author reading and two independent
whole-proof readings have passed.
The result concerns a bounded map from the actual rational local
Hilbert completion to the actual global closed theta ideal.
The smooth space called \(E\) on Masullo p.~515 is kept distinct
from this completion.  The cutoff hypotheses are indispensable:
bounded translation-equivariance alone is not excluded.

Fix \(S=\{\infty\}\) over \(\mathbb Q\), the printed twist
\(A=1/2\), and the single-line norm \(I=\{0\}\).  Put
\[
 \begin{gathered}
 p_s(x)=(1+x^2)^{s/2},\qquad
 H_s^{\rm pr}=L^2(\mathbb R,p_s(x)\,dx),\qquad s\ge0,\\
 \mathscr X_0=\{f\in\mathcal S(\mathbb R)^{\rm even}:
                       f(0)=0,\ \int_{\mathbb R}f(t)\,dt=0\},\\
 \mathscr X_{\rm fl}=\mathscr X_0\cap
                         \mathcal S_{00}(\mathbb R)^{\rm even}.
 \end{gathered}
 \tag{MA-PI.210}
\]
Here \(\mathcal S_{00}(\mathbb R)\) is the actual Schwartz
space flat to every order at the additive origin.  In particular,
an even Schwartz function vanishing in a neighborhood of zero
belongs to it.  The topology on \(\mathscr X_0\) is the induced
Schwartz Fr\'echet topology.

The local chart, including its scalar, is
\[
 (Vf)(x)=2e^{x/2}f(e^x),\qquad
 \tau_yh(x)=h(x-y),\qquad
 (D_yf)(t)=e^{-y/2}f(e^{-y}t).
 \tag{MA-PI.211}
\]
Let \(\mathscr I\) be the original rational centered smooth
global theta ideal and \(J_s^{\rm pr}=\overline{\mathscr I}^{H_s^{\rm pr}}\).
For the obstruction below assume \(s>1\), and let \(P_Y\)
denote the exact printed global orthogonal PW projector at
physical level \(Y\), transported by the unitary (MA-PI.188).
The local chart and density statements hold already for every
\(s\ge0\).

For each \(Y\ge Y_0>0\), fix a real smooth ball cutoff
\(\rho_Y\) with \(0\le\rho_Y\le1\), equal to one on
\([0,Y/2]\) and to zero on \([2Y,\infty)\), and set
\[
 B_Y=\mathcal F_\infty^{-1}M_{\rho_Y(|\xi|)}\mathcal F_\infty,
 \qquad
 (\mathcal F_\infty f)(\xi)=
                  \int_{\mathbb R}f(t)e^{-2\pi it\xi}\,dt.
 \tag{MA-PI.212}
\]
This defines a continuous map on even Schwartz space for each
fixed \(Y\).  Its output need not satisfy either the flatness
condition or the value-at-zero condition in (MA-PI.210).

\paragraph{MA-PI-010 (no bounded local-to-global transport with these cutoff identities).}
The completion of the source's smooth local space
\[
 E=|\cdot|^{1/2}H_{0,\rm cont}
                (\{\pm1\},\mathcal S_{00}(\mathbb R))
 \tag{MA-PI.213}
\]
in the exact printed \(A=1/2,I=\{0\}\) norm is isometrically
\(H_s^{\rm pr}\) through \(V\).  The separately printed local
two-moment \(S_0\) completion has the same isometric realization.
The images \(V\mathscr X_0\) and \(V\mathscr X_{\rm fl}\) are
both dense in this Hilbert space.

Let \(s>1\).  Suppose
\(R:H_s^{\rm pr}\to J_s^{\rm pr}\) is a bounded linear map
such that \(R\tau_y=\tau_yR\) for every real \(y\).
Assume, for every \(Y\ge Y_0\), the following plateau
identities on the \emph{full} domain \(\mathscr X_0\):
\[
 \begin{array}{ll}
 P_YRVf=RVf,& B_Yf=f,\quad f\in\mathscr X_0,\\
 P_YRVf=0,& B_Yf=0,\quad f\in\mathscr X_0.
 \end{array}
 \tag{MA-PI.214}
\]
Then \(R=0\).

The same conclusion follows from the stronger, precisely
defined equation only on the flat two-moment core:
\[
       P_YRVf=RV(B_Yf),
       \qquad f\in\mathscr X_{\rm fl},\quad Y\ge Y_0.
 \tag{MA-PI.215}
\]
The right side of (MA-PI.215) uses the continuous extension
of \(V\) from all even Schwartz functions into the local
completion.  It does not assert \(B_Yf\in\mathscr X_0\).
Plateau identities merely on \(\mathscr X_{\rm fl}\), without
the full equation (MA-PI.215), are not a hypothesis of this
second conclusion.

\begin{proof}
\emph{1. The actual quotient, measure, and local completion.}
For the finite sign group the separated coinvariant quotient
is exactly the even subspace: the odd part of \(f\) is
\(\tfrac12(f-f(-\cdot))\), a relation, and the even projection
is a continuous splitting.  Thus no unfinished general
S-unit topological identification is needed here.
Definition~2.3.1.1 and (2.3.1.2), pp.~99--100, give
\(\Phi_{\mathbb Q,\{\infty\}}f(t)=f(t)+f(-t)\).
For even \(f\), its half-centered logarithmic realization
is exactly (MA-PI.211).

The additive Haar measure is \(dt\), and the multiplicative
measure is \(dt/|t|\), by Definitions~A.3.1.2(1) and
A.3.2.2(1), pp.~556 and~560.  The sign group has counting
measure by Definition~A.3.4.1(3)(2), pp.~561--562.
The two-sign Weil formula therefore gives the quotient
measure \(dt/t\) on \(\mathbb R_{>0}\), hence \(dx\) on
the logarithmic line.  Its norm-one fiber has mass one.
The norm (10.1.1.5), p.~434, is applied to the uncentered
local theta function; equivalently, the already half-centered
element of \(E\) has the printed \(p_s\) norm.  There is
no second half twist.  The exact resulting identity is
\[
 \begin{aligned}
 \|Vf\|_{H_s^{\rm pr}}^2
   &=4\int_0^\infty |f(t)|^2p_s(\log t)\,dt\\
   &=2\int_{\mathbb R}|f(t)|^2p_s(\log|t|)\,dt
                   \qquad(f\text{ even}).
 \end{aligned}
 \tag{MA-PI.216}
\]
The factor is fixed; it is not an unknown comparison scalar.
Because \(p_s(\log t)\) is integrable near zero and grows
more slowly than any positive power of \(t\) at infinity,
(MA-PI.216) proves that
\(V:\mathcal S(\mathbb R)^{\rm even}\to H_s^{\rm pr}\)
is continuous.  For example a sufficiently large zeroth
Schwartz seminorm gives an integrable majorant at both ends.

For \(v\in C_c^\infty(\mathbb R)\), the formula
\[
 f_v(t)=\tfrac12|t|^{-1/2}v(\log|t|)\quad(t\ne0),
                  \qquad f_v(0)=0
 \tag{MA-PI.217}
\]
defines an even smooth function compactly supported away
from zero, and \(Vf_v=v\).  Hence the image of the flat
smooth space in (MA-PI.213) contains every compact smooth
logarithmic function.  These are dense in \(H_s^{\rm pr}\).
Since (MA-PI.216) is the exact defining norm, its completion
is precisely that Hilbert space, by an isometry.
This statement completes the smooth locally convex space
\(E\); it does not change the topology assigned to \(E\)
itself on pp.~515--516.

\emph{2. Density with both principal conditions retained.}
For the vector in (MA-PI.217), the remaining condition is
\[
 \int_{\mathbb R}f_v(t)\,dt
     =\ell_+(v),\qquad
 \ell_+(v)=\int_{\mathbb R}e^{x/2}v(x)\,dx.
 \tag{MA-PI.218}
\]
This is the plus-half moment.  It differs from the minus-half
moment of the physical \emph{frequency} chart in (MA-PI.197).
Choose \(\beta\in C_c^\infty\) with \(\ell_+(\beta)=1\).
For \(L\to+\infty\), put
\[
 \begin{gathered}
 \beta_L(x)=e^{-L/2}\beta(x-L),\qquad
 \ell_+(\beta_L)=1,\\
 \|\beta_L\|_{H_s^{\rm pr}}
       \le C_{s,\beta}e^{-L/2}(1+L)^{s/2}\longrightarrow0,\\
 v_L=v-\ell_+(v)\beta_L\longrightarrow v
                                      \quad\text{in }H_s^{\rm pr}.
 \end{gathered}
 \tag{MA-PI.219}
\]
The inverse image of \(v_L\) in (MA-PI.217) is a genuine
even compact function away from zero, with integral zero.
It therefore lies in \(\mathscr X_{\rm fl}\), including its
value-at-zero condition.  Density of compact smooth \(v\)
now proves
\[
 \overline{V\mathscr X_{\rm fl}}^{\,H_s^{\rm pr}}
   =\overline{V\mathscr X_0}^{\,H_s^{\rm pr}}
   =H_s^{\rm pr}.
 \tag{MA-PI.220}
\]
The source's definition \(S_0=\ker\delta_0\cap\ker\widehat\delta_0\),
(7.2.2.3), p.~340, is exactly \(\mathscr X_0\) in the even
rational sector.  Thus its printed local completion has
the same isometric closure.  No continuity of \(\ell_+\)
on the whole Hilbert completion is asserted or used.

\emph{3. Normalized equivariance and the full-domain obstruction.}
Direct substitution gives
\[
       VD_yf=\tau_yVf.
 \tag{MA-PI.221}
\]
For the raw dilation \(L_yf(t)=f(e^{-y}t)\), the equation
instead is \(VL_y=e^{y/2}\tau_yV\).  Thus (MA-PI.221)
specifies exactly the normalized action being intertwined,
in accordance with the source's warning in Remark~10.1.1.2,
pp.~434--435.  The domain \(\mathscr X_0\) is invariant
under these dilations.

Set \(T=RV|_{\mathscr X_0}\).  By (MA-PI.216), it is a
continuous linear map from the full two-moment Schwartz
space into \(J_s^{\rm pr}\); by (MA-PI.221),
\(TD_y=\tau_yT\).  Hypothesis (MA-PI.214) is exactly the
principal Hilbert-target instance of (MA-PI.195).
The complete full-Schwartz theorem MA-PI-009, using
MA-PI-008 for this \(s>1\) target, gives \(T=0\).
Equation (MA-PI.220) and boundedness of \(R\) then imply
\(R=0\) on the entire completed local space.
This invocation retains the known critical-line simple-zero
count from Conrey (1989), with its exact statement and
counting scope recorded in MA-PI-008; it does not assume
\(\mathsf A\), \(M\), or the Riemann hypothesis.

\emph{4. The fixed-level graph approximation for (MA-PI.215).}
It remains to justify why that stronger equation on the
smaller core implies the required full-domain plateaus.
Choose an even \(\chi\in C_c^\infty\), equal to one near
zero, and a fixed even annular \(b\in C_c^\infty\) with
\(\int b=1\).  For \(f\in\mathscr X_0\), set
\[
 \begin{gathered}
 r_\varepsilon(t)=\chi(t/\varepsilon)f(t),\qquad
 a_\varepsilon=\int_{\mathbb R}r_\varepsilon(t)\,dt,\\
 f_\varepsilon=f-r_\varepsilon+a_\varepsilon b.
 \end{gathered}
 \tag{MA-PI.222}
\]
For sufficiently small \(\varepsilon\), this is even,
Schwartz, zero near the origin and of integral zero.
Thus \(f_\varepsilon\in\mathscr X_{\rm fl}\).
Evenness and \(f(0)=0\) give \(f(t)=O_f(t^2)\) near zero.
Consequently
\[
 \begin{aligned}
 |a_\varepsilon|&\le C_f\varepsilon^3,\\
 \|Vr_\varepsilon\|_{H_s^{\rm pr}}^2
     &\le C_{f,s}\varepsilon^5
                         (1+|\log\varepsilon|)^s,\\
 Vf_\varepsilon&\longrightarrow Vf
                                  \quad\text{in }H_s^{\rm pr}.
 \end{aligned}
 \tag{MA-PI.223}
\]
The middle estimate follows from (MA-PI.216) and the
integral of \(t^4p_s(\log t)\) on \((0,C\varepsilon)\).

Now fix \(Y\), without imposing any estimate uniform in
\(Y\).  Differentiation under the Fourier integral gives,
for each integer \(k\ge0\),
\[
 \sup_{\xi\in\mathbb R}
     |\partial_\xi^k\widehat r_\varepsilon(\xi)|
       \le (2\pi)^k\int |t|^k|r_\varepsilon(t)|\,dt
       \le C_{f,k}\varepsilon^{k+3}.
 \tag{MA-PI.224}
\]
The fixed multiplier \(\rho_Y(|\xi|)\) is smooth and
compactly supported, so (MA-PI.224) implies
\(\rho_Y\widehat r_\varepsilon\to0\) in the fixed-support
\(C_c^\infty\) topology.  Fourier inversion is continuous
from that space into Schwartz space.  Since
\(a_\varepsilon\to0\), it follows that
\[
 B_Yf_\varepsilon\longrightarrow B_Yf
            \quad\text{in }\mathcal S(\mathbb R)^{\rm even},
 \qquad
 V(B_Yf_\varepsilon)\longrightarrow V(B_Yf)
            \quad\text{in }H_s^{\rm pr}.
 \tag{MA-PI.225}
\]
Thus the approximants converge in exactly the two graph
coordinates needed for (MA-PI.215).  They need not converge
to \(f\) in the full Schwartz topology before applying
\(B_Y\).  No bounded realization of \(B_Y\) on the whole
weighted local Hilbert space has been assumed.

Apply (MA-PI.215) to \(f_\varepsilon\).  Boundedness of
\(R\) and \(P_Y\), together with (MA-PI.223)--(MA-PI.225),
gives
\[
       P_YRVf=RV(B_Yf),\qquad f\in\mathscr X_0.
 \tag{MA-PI.226}
\]
For a plateau input this is (MA-PI.214), and Step~3 proves
\(R=0\).  Every term on the right is defined because \(V\)
acts on all even Schwartz functions, including outputs
\(B_Yf\) with nonzero value at the origin.
\end{proof}

\paragraph{Why the proof does not transfer flat-core plateau hypotheses.}
If \(f\in\mathcal S_{00}(\mathbb R)\) satisfies \(B_Yf=f\),
then \(\widehat f\) is supported in \([-2Y,2Y]\).
Its inverse Fourier integral extends \(f\) to an entire
function.  All derivatives at zero vanish, so this entire
function is identically zero.  In particular, the inner
plateau on the flat core supplies no nonzero test pair
for the separated-support argument of MA-PI-008.
The full-domain plateau assumption (MA-PI.214) and the
stronger flat-core equation (MA-PI.215) must therefore
remain distinct.  The graph approximation just proved
transfers the latter equation; it does not transfer
plateau membership of its approximants.
This observation does not construct a nonzero map satisfying
only the flat-core plateau constraints or decide whether
their outside-plateau part alone forces vanishing.

\paragraph{Exact source and remaining scope.}
The preceding source identifications use Masullo,
Definition~2.3.1.1, pp.~99--100; the flat-space definition
5.3.4.1, p.~259, also stated in Theorem~5.3.2.1(1), p.~243;
the two-moment definition (7.2.2.3), p.~340;
Definition~9.1.2.1 and (9.1.2.3), pp.~399--400; and
Definition~10.1.1.1(2), (10.1.1.5), p.~434.
Page~515 defines the smooth \(E\) in (MA-PI.213), whereas
the local \(W_{1/2}^{s,2}\) completion is a separate object
displayed above it.  Page~516 explicitly describes the
smooth source spaces as Hausdorff locally convex spaces.
The physical ball is p.~516, Step~II, and (12.4.1.8),
p.~517, uses Fourier multiplication followed by inverse
Fourier transform.  The eventual plateau conventions are
those on p.~519.

For the global target, the compact-fiber Haar mass is
retained by the exact unitary (MA-PI.188), rather than
discarded by a comparison of equivalent norms.
A consistent nonzero rescaling of the local chart or
positive rescaling of its norm would preserve density
and the zero conclusion, but no such rescaling is needed
in (MA-PI.216).  An arbitrary equivalent norm would not
license transport of an orthogonal projector.

The theorem excludes a fixed bounded local-to-global map
with the stated ideal range, normalized equivariance and
all-large-level cutoff identities.  It does not exclude
bounded equivariant maps without those identities,
construct an inverse for a global-to-local comparison,
or rule out level-dependent, altered-cutoff, or
Gram-corrected constructions.  It makes no Hilbert-target
claim for \(s\le1\).  The broader scalar arithmetic identity
in MA-FT-002 remains a separate question.

In the inherited cote~1 BR-TR-01/02 placement,
scans079--080 and111--112, this is an obstruction to an
actual bounded comparison before nuclear trace cyclicity.
The finite sign coinvariant is an explicit split quotient;
the norm-half centering in (MA-PI.211) is analytic and is
not an integral Tate twist.  No finite-length \(K_0\),
weight, or purity assertion supplies the missing cutoff
intertwiner or a new arithmetic positivity theorem.
